\documentclass[11pt,a4paper]{article}

\usepackage[utf8]{inputenc}
\usepackage[T1]{fontenc}
\usepackage[spanish,es-nodecimaldot]{babel}
\usepackage{amsmath,amssymb,bm}
\usepackage{booktabs}
\usepackage{array}
\usepackage{graphicx}
\usepackage{geometry}
\usepackage{xcolor}
\usepackage{hyperref}
\usepackage{siunitx}

\geometry{margin=2.2cm}
\setlength{\parindent}{0pt}
\setlength{\parskip}{0.55em}

\newcommand{\dd}{\mathrm{d}}
\newcommand{\R}{\mathcal{R}}

\title{\textbf{Avance de resultados: régimen phantom en gravedad unimodular con difusión}}
\author{}
\date{}

\begin{document}
\maketitle

\section{Modelo y ecuaciones de evolución}

Se considera un universo FLRW espacialmente plano, homogéneo e isotrópico,
con un fluido barotrópico y una función de difusión de la forma
\begin{equation}
    p=\omega\rho,
    \qquad
    Q=\alpha\rho^\beta.
\end{equation}
En las unidades empleadas en el proyecto y tomando nula la constante de
integración cosmológica explícita, las ecuaciones de fondo son
\begin{align}
    3H^2 &= \rho+\alpha\rho^\beta
    =\rho F(\rho),\\
    F(\rho)&=1+\alpha\rho^{\beta-1},\\
    D(\rho)&=1+\alpha\beta\rho^{\beta-1},\\
    \dot{\rho}
    &=
    -\frac{3H(1+\omega)\rho}{D(\rho)},\\
    \dot H&=-\frac{1+\omega}{2}\rho.
\end{align}

Para el régimen phantom se define
\begin{equation}
    \gamma\equiv-(1+\omega)>0,
\end{equation}
de modo que
\begin{equation}
    \dot\rho=
    \frac{3\gamma H\rho}{D(\rho)}.
\end{equation}
Por lo tanto, en una rama expansiva el signo de $D(\rho)$ controla
directamente si la densidad aumenta o disminuye.

\section{Barrido numérico completo}

El barrido \texttt{full} contiene $5425$ puntos en total. El sector phantom
corresponde a
\[
\omega=-1.4,\quad -1.2,\quad -1.05,
\]
con $775$ puntos por cada valor de $\omega$, es decir, $2325$ puntos phantom.

No se registraron errores de integración en el archivo combinado.
La frontera de tiempo finito/infinito prevista analíticamente en
$\beta=2$ fue contrastada en $1386$ puntos del sector
$\alpha>0$, $\rho\rightarrow\infty$, obteniéndose acuerdo numérico completo
con la predicción asintótica.

\section{Clasificación final del sector phantom}

Después de combinar el barrido numérico con las relaciones asintóticas,
la clasificación queda:

\begin{table}[ht]
\centering
\small
\begin{tabular}{p{4.0cm} p{3.0cm} p{6.0cm}}
\toprule
\textbf{Región} & \textbf{Destino} & \textbf{Comportamiento futuro}\\
\midrule
$\beta<0$, $\alpha>0$, $D_0<0$
&
Rip asintótico dominado por difusión
&
$t\to\infty$, $\rho,p\to0$, mientras
$Q,a,H,\R\to\infty$ y $\dot H\to0^{+}$.
\\[1mm]

$\alpha<0$, $\beta>1$
&
Tipo IV en $D(\rho_D)=0$
&
$t_s<\infty$; $\rho,p,a,H,\dot H$ finitos,
pero $\dot\rho$ y $\ddot H$ divergen.
\\[1mm]

$\alpha>0$, $\beta\ge2$
&
Little Rip
&
$t\to\infty$; $\rho,|p|,Q,a,H,\dot H\to\infty$.
No existe singularidad a tiempo cósmico finito.
\\[1mm]

Resto de las ramas phantom físicas
&
Big Rip, Tipo I
&
$t_s<\infty$; $a,\rho,|p|,H\to\infty$.
\\
\bottomrule
\end{tabular}
\caption{Clasificación analítico-numérica del régimen phantom.}
\end{table}

Para $\rho_0=1$, la condición de la primera región se reduce a
\begin{equation}
    D_0=1+\alpha\beta<0
    \quad\Longleftrightarrow\quad
    \alpha>-\frac{1}{\beta},
    \qquad \beta<0.
\end{equation}

Cada slice phantom del barrido contiene, tras la revisión analítica,
\begin{align*}
417&\quad \text{Big Rip},\\
158&\quad \text{rips dominados por difusión},\\
140&\quad \text{Little Rip},\\
60&\quad \text{Tipo IV}.
\end{align*}
La igualdad de estos conteos para los tres valores de $\omega<-1$ estudiados
indica que las fronteras cualitativas dependen de $(\alpha,\beta)$, mientras
que $|1+\omega|$ fija principalmente la escala temporal. Los porcentajes
del barrido no deben interpretarse como probabilidades físicas, ya que
dependen de la discretización escogida.

\begin{figure}[ht]
\centering
\IfFileExists{mapa_phantom_omega_m1p2_revisado.png}{
\includegraphics[width=0.88\textwidth]{mapa_phantom_omega_m1p2_revisado.png}
}{
\fbox{\parbox{0.8\textwidth}{Subir junto al .tex la figura
\texttt{mapa\_phantom\_omega\_m1p2\_revisado.png}.}}
}
\caption{Mapa revisado de destinos cosmológicos para $\omega=-1.2$.}
\end{figure}

\section{Frontera analítica $\beta=2$}

Para $\alpha>0$, $\beta>1$ y $\rho\to\infty$,
\begin{equation}
    H\sim\sqrt{\frac{\alpha}{3}}\rho^{\beta/2},
    \qquad
    D\sim\alpha\beta\rho^{\beta-1}.
\end{equation}
Por consiguiente,
\begin{equation}
    \dot\rho
    \sim
    \frac{\sqrt3\,\gamma}{\beta\sqrt{\alpha}}
    \rho^{\,2-\beta/2}.
\end{equation}
El tiempo restante satisface
\begin{equation}
    \Delta t
    \propto
    \int^\infty
    \rho^{-(2-\beta/2)}\,\dd\rho.
\end{equation}
La integral converge si y sólo si $\beta<2$. Por tanto,
\begin{equation}
    \boxed{
    1<\beta<2 \Rightarrow \text{Big Rip},
    \qquad
    \beta\ge2 \Rightarrow \text{Little Rip}
    }
\end{equation}
para $\alpha>0$. El caso $\beta=2$ es marginal: $\rho$ crece
exponencialmente con el tiempo y la divergencia se desplaza a
$t\to\infty$.

\section{Frontera $D(\rho)=0$ y singularidad Tipo IV}

Para $\alpha<0$ y $\beta>1$ existe
\begin{equation}
    \rho_D=
    \left(
    -\frac{1}{\alpha\beta}
    \right)^{1/(\beta-1)}.
\end{equation}
En esta densidad,
\begin{equation}
    F(\rho_D)
    =
    1-\frac{1}{\beta}
    =
    \frac{\beta-1}{\beta}>0,
\end{equation}
por lo que $H_D$ permanece real y finito. También permanecen finitos
$\rho_D$, $p_D$, $a_D$ y $\dot H_D$. Sin embargo,
\[
D(\rho)\to0
\]
hace divergir $\dot\rho$, y como
\begin{equation}
    \ddot H=\frac{\gamma}{2}\dot\rho,
\end{equation}
también diverge $\ddot H$. Cerca de la frontera,
\begin{equation}
    |\rho-\rho_D|
    \propto
    (t_s-t)^{1/2},
\end{equation}
de donde
\begin{equation}
    |\dot\rho|,
    |\ddot H|
    \propto
    (t_s-t)^{-1/2}.
\end{equation}
Esto confirma la clasificación Tipo IV.

\section{Caso F2: $\omega=-1.2$, $\alpha=4$, $\beta=-1$}

Para este caso,
\begin{equation}
    Q=\frac{4}{\rho},
    \qquad
    D(\rho)=1-\frac{4}{\rho^2}.
\end{equation}
Con $\rho_0=1$,
\[
D_0=-3<0,
\]
por lo que, a pesar de que el fluido es phantom,
\begin{equation}
    \boxed{\dot\rho<0}.
\end{equation}
La densidad material disminuye hacia cero, pero
\begin{equation}
    Q=\frac4\rho\to\infty,
\end{equation}
y la ecuación de Friedmann produce
\begin{equation}
    3H^2=\rho+\frac4\rho
    \quad\Longrightarrow\quad
    H\to\infty.
\end{equation}

\subsection{Asintótica general de la rama F2-like}

Para $\alpha>0$, $\beta<0$ y $D<0$, cuando $\rho\to0$,
\begin{align}
    H&\sim
    \sqrt{\frac{\alpha}{3}}\rho^{\beta/2},\\
    \dot\rho
    &\sim
    -\frac{\sqrt3\,\gamma}
    {|\beta|\sqrt{\alpha}}
    \rho^{\,2-\beta/2}.
\end{align}
Integrando,
\begin{equation}
    \rho(t)\propto
    t^{-2/(2-\beta)}.
\end{equation}
Entonces
\begin{align}
    Q(t)&\propto
    t^{-2\beta/(2-\beta)}
    \longrightarrow\infty,\\
    H(t)&\propto
    t^{-\beta/(2-\beta)}
    \longrightarrow\infty,\\
    \dot H(t)&\propto
    t^{-2/(2-\beta)}
    \longrightarrow0^+.
\end{align}
Además,
\begin{equation}
    \ln a(t)
    \propto
    t^{\,2(1-\beta)/(2-\beta)},
\end{equation}
por lo que $a(t)$ crece más rápido que exponencialmente.

\subsection{Relación exacta $a(\rho)$}

Para $\beta\neq1$,
\begin{equation}
    \ln\frac{a}{a_0}
    =
    \frac{1}{3\gamma}
    \left[
    \ln\frac{\rho}{\rho_0}
    +
    \frac{\alpha\beta}{\beta-1}
    \left(
    \rho^{\beta-1}
    -
    \rho_0^{\beta-1}
    \right)
    \right].
\end{equation}
En F2, con $a_0=\rho_0=1$,
\begin{equation}
    \boxed{
    a(\rho)=
    \rho^{5/3}
    \exp\left[
    \frac{10}{3}
    \left(
    \rho^{-2}-1
    \right)
    \right]
    }.
\end{equation}
De aquí se demuestra directamente que
\[
\rho\to0
\quad\Longrightarrow\quad
a\to\infty.
\]

Para F2 se obtiene en particular
\begin{align}
    \rho(t)&\sim(0.2598\,t)^{-2/3},\\
    H(t)&\sim0.737\,t^{1/3},\\
    \dot H(t)&\sim0.246\,t^{-2/3},\\
    \ln a(t)&\sim0.553\,t^{4/3}.
\end{align}

\section{Interpretación geométrica de F2}

Definiendo cantidades efectivas
\begin{equation}
    \rho_{\rm eff}=\rho+Q,
    \qquad
    p_{\rm eff}=p-Q,
\end{equation}
se tiene
\begin{equation}
    w_{\rm eff}
    =
    \frac{\omega\rho-Q}{\rho+Q},
\end{equation}
y
\begin{equation}
    w_{\rm eff}+1
    =
    \frac{(1+\omega)\rho}{\rho+Q}.
\end{equation}
Por tanto, en la rama F2,
\begin{equation}
    w_{\rm eff}\to-1^{-},
\end{equation}
pero $H$ no tiende a una constante: diverge. En consecuencia,
la solución no es asintóticamente de Sitter.

El escalar de Ricci para FLRW plano es
\begin{equation}
    \R=6(\dot H+2H^2).
\end{equation}
Como $H\to\infty$,
\[
\R\to\infty
\qquad (t\to\infty).
\]
También
\[
\frac{\ddot a}{a}=H^2+\dot H\to\infty.
\]
Por ello el régimen es genuinamente ``rip-like'', aunque no presenta una
singularidad a tiempo cósmico finito.

\section{Estado de la interpretación de F2}

El nombre \texttt{diffusion\_driven\_infinite\_time\_rip} debe entenderse
como una etiqueta de trabajo, no como una clasificación establecida en la
literatura. Geométricamente comparte características con los escenarios
Little Rip y Little Sibling of the Big Rip: $a,H$ divergen a tiempo
infinito y la curvatura crece sin cota. Su rasgo distintivo es que la
densidad del fluido fundamental satisface
\[
\rho\to0,
\]
mientras que es la función de difusión la que diverge:
\[
Q\to\infty.
\]
Por ahora debe presentarse como \emph{candidato a régimen asintótico
dominado por difusión}, no como una nueva clase de singularidad demostrada.

\section{Advertencia termodinámica}

La relación de Gibbs combinada con la ecuación de no conservación de UG
conduce, bajo las hipótesis termodinámicas usuales, a
\begin{equation}
    T\,\dd S=-V\,\dd Q.
\end{equation}
En toda la rama F2-like,
\[
\alpha>0,\qquad
\beta<0,\qquad
\dot\rho<0,
\]
por lo que
\begin{equation}
    \dot Q
    =
    \alpha\beta\rho^{\beta-1}\dot\rho
    >0.
\end{equation}
Si $T>0$ y $V>0$, esto implica $\dot S<0$. Por tanto, el régimen es
matemáticamente consistente como solución de las ecuaciones de fondo,
pero presenta una tensión con la segunda ley bajo la formulación
termodinámica estándar. Esta cuestión debe estudiarse antes de afirmar
viabilidad física.

\section{Conclusiones para presentar}

\begin{enumerate}
    \item El barrido completo revela cuatro destinos robustos en el sector
    phantom: Big Rip, Tipo IV, Little Rip y un régimen asintótico dominado
    por difusión.
    \item La frontera $\beta=2$ separa Big Rip y Little Rip para
    $\alpha>0$ y fue confirmada numéricamente en todos los puntos
    comparables del barrido.
    \item Para $\alpha<0$, $\beta>1$, la evolución alcanza $D=0$ a tiempo
    finito y genera una singularidad Tipo IV.
    \item Para $\beta<0$, $\alpha>0$ y $D_0<0$, la densidad phantom
    disminuye hasta cero pero $Q$, $H$, $a$ y la curvatura divergen a
    tiempo infinito.
    \item F2 $(\omega,\alpha,\beta)=(-1.2,4,-1)$ es un representante limpio
    de esta última región y admite una caracterización asintótica y una
    relación exacta $a(\rho)$.
    \item La posible novedad de F2 reside en el mecanismo
    $\rho\to0$ con $Q\to\infty$, no todavía en haber demostrado una nueva
    clase universal de singularidad.
\end{enumerate}

\begin{thebibliography}{9}

\bibitem{Corral2020}
C.~Corral, N.~Cruz y E.~González,
\emph{Diffusion in unimodular gravity: Analytical solutions, late-time
acceleration, and cosmological constraints},
Phys.\ Rev.\ D \textbf{102}, 023508 (2020).

\bibitem{ThermoUG2024}
M.~Cruz, N.~Cruz y S.~Lepe,
\emph{Exploring thermodynamics inconsistencies in unimodular gravity:
a comparative study of two energy diffusion functions},
Eur.\ Phys.\ J.\ C \textbf{84}, 1186 (2024).

\bibitem{LittleRip}
P.~H.~Frampton, K.~J.~Ludwick, S.~Nojiri, S.~D.~Odintsov y R.~J.~Scherrer,
\emph{Models for Little Rip Dark Energy},
arXiv:1108.0067.

\bibitem{LSBR}
M.~Bouhmadi-López, A.~Errahmani, P.~Martín-Moruno, T.~Ouali y Y.~Tavakoli,
\emph{The Little Sibling of the Big Rip singularity},
arXiv:1407.2446.

\end{thebibliography}

\end{document}
